\begin{frame}{Ecuación diferencial radial}

    Considerando el potencial coulombiano
    \[
        V(r)=-\frac{e^2}{4\pi\epsilon_0r},
    \]
    y la solución angular obtenida anteriormente,
    $\lambda=l(l+1)$, la ecuación radial resulta

    \begin{equation}
    -\frac{\hbar^2}{2mr^2}
    \frac{d}{dr}\left(r^2\frac{dR}{dr}\right)
    +
    \left[
        \frac{\hbar^2l(l+1)}{2mr^2}
        -\frac{e^2}{4\pi\epsilon_0r}
    \right]R=ER
    \label{eq:ec de schrodinger radial}
    \end{equation}


    El término proporcional a $l(l+1)$ representa
    la contribución centrífuga.

\end{frame}

\begin{frame}{Transformación de la ecuación radial}

    Definimos una nueva función
    \begin{equation}
        u(r)\equiv rR(r)
        \label{eq: definición u}
    \end{equation}
        

    De esta definición obtenemos
    \[
        R'=\frac{ru'-u}{r^2},
        \qquad
        \frac{d}{dr}(r^2R')=ru''.
    \]

    Sustituyendo en la ecuación radial:
    \[
    -\frac{\hbar^2}{2m}u''
    +
    \left[
        \frac{\hbar^2l(l+1)}{2mr^2}
        -\frac{e^2}{4\pi\epsilon_0r}
    \right]u=Eu.
    \]

\end{frame}

\begin{frame}{Adimensionalización de la ecuación}

    Para estados ligados tenemos $E<0$. Definimos
    \begin{equation}
        \rho=\frac{\sqrt{2m|E|}}{\hbar}r
    \end{equation}

    y el parámetro
    \begin{equation}
        \rho_0=\frac{me^2}{2\pi\epsilon_0\hbar\sqrt{2m|E|}}
    \end{equation}
 
    De esta manera, la ecuación radial queda
    \begin{equation}
        \frac{d^2u}{d\rho^2} =
        \left[
            1-\frac{\rho_0}{\rho}+\frac{l(l+1)}{\rho^2}
        \right]u
        \label{eq:edo radial adimensional}
    \end{equation}

    La dependencia energética se encuentra ahora contenida
    en el parámetro $\rho_0$.

\end{frame}

\begin{frame}{Comportamiento asintótico de la solución}

    \begin{columns}[T]

        \begin{column}{0.48\textwidth}
            \textbf{Para $\rho\rightarrow\infty$}

            La ecuación se aproxima a
            \[
                u''\approx u.
            \]

            Su solución general es
            \[
                u=Ae^{-\rho}+Be^\rho.
            \]

            La normalización exige $B=0$. Por lo tanto,
            \[
                u\sim Ae^{-\rho}.
            \]
        \end{column}

        \begin{column}{0.48\textwidth}
            \textbf{Para $\rho\rightarrow0$}

            La ecuación se aproxima a
            \[
                u''\approx
                \frac{l(l+1)}{\rho^2}u.
            \]

            Proponiendo $u\sim\rho^s$:
            \[
                s(s-1)=l(l+1).
            \]

            Sus raíces son $s=l+1$ y $s=-l$
            La regularidad en el origen selecciona
            \[
                u\sim B\rho^{l+1}.
            \]
        \end{column}
    \end{columns}
\end{frame}

\begin{frame}{Propuesta de solución}

    A partir del comportamiento asintótico, proponemos
    \[
        u(\rho)=\rho^{l+1}e^{-\rho}v(\rho)
    \]

    Derivando dos veces:
    \begin{equation}
        u''={}\rho^{l+1}e^{-\rho}v''+2\left(\frac{l+1}{\rho}-1\right)v'+
        \left(
            \frac{l(l+1)}{\rho^2}
            -\frac{2(l+1)}{\rho}+1
        \right)v
        \label{eq:dud2}
    \end{equation}

    Sustituyendo la ecuación \ref{eq:dud2} en la ecuación diferencial radial adimensional \ref{eq:edo radial adimensional}:
    
    \begin{equation}
        \rho v''+2(l+1-\rho)v'+[\rho_0-2(l+1)]v=0
    \end{equation}
\end{frame}

\begin{frame}{Desarrollo en series de potencias}

    Proponemos una solución de la forma
    \[
        v(\rho)=\sum_{k=0}^{\infty}c_k\rho^k.
    \]

    Sus derivadas son
    \[
        v'=\sum_{k=0}^{\infty}(k+1)c_{k+1}\rho^k,
    \]
    \[
        v''=\sum_{k=0}^{\infty}k(k+1)c_{k+1}\rho^{k-1}.
    \]

\end{frame}
\begin{frame}{Desarrollo en series de potencias}
    Sustituyendo en la ecuación diferencial y agrupando
    los coeficientes de cada potencia de $\rho$:
    \[
    \begin{aligned}
    (k+1)(k+2l+2)c_{k+1}+[\rho_0-2(k+l+1)]c_k=0.
    \end{aligned}
    \]

    Obtenemos la relación de recurrencia
    \begin{equation}
        c_{k+1}=
        \frac{2(k+l+1)-\rho_0}
        {(k+1)(k+2l+2)}c_k
    \end{equation}
        
\end{frame}
\begin{frame}{Terminación de la serie y cuantización}

    Para obtener una solución polinómica de grado $N$,
    exigimos $c_{N-1}\neq0$ y  $c_{N+1}=0.$


    Aplicando esta condición a la relación de recurrencia:
    \[
        2(N+l+1)-\rho_0=0.
    \]

        Definimos el número cuántico principal
    \[
        n=N+l+1
    \]

    De esta manera:
    \[
        \rho_0=2n
    \]

    y, como $N\geq0$ y $l\geq0$:
    \[
        n=1,2,\ldots,\qquad 0\leq l\leq n-1.
    \]

\end{frame}
\begin{frame}{Energías permitidas}
    Como $\rho_0=2n$, despejando la energía:
    \[
        E_n=
        -\frac{me^4}{2\hbar^2(4\pi\epsilon_0)^2}
        \frac{1}{n^2} \simeq-\frac{13.6\,\text{eV}}{n^2}\qquad n=1,2,\ldots
    \]
    Definiendo el radio de Bohr
    \[
        a_0=\frac{4\pi\epsilon_0\hbar^2}{me^2},
    \]
    obtenemos:
    \[
        \rho=\frac{r}{na_0}.
    \]
\end{frame}
\begin{frame}{Solución radial y función de onda}
    Recordando que $u=rR$ y
    \(
        u(\rho)=\rho^{l+1}e^{-\rho}v(\rho),
    \)
    obtenemos
    \begin{equation}
        R_{nl}(r)=
        C_{nl}\frac{\rho^{l+1}e^{-\rho}}{r}
        v_{nl}(\rho)
    \end{equation}
        
    La función de onda completa queda expresada como
    \[
        \boxed{
        \psi_{nlm}(\mathbf r)=
        C_{nl}\sqrt{\frac{2l+1}{4\pi}\frac{(l-m)!}{(l+m)!}}\frac{\rho^{l+1}e^{-\rho}}{r}v_{nl}(\rho)\,P_l^m(\cos\theta)\, e^{im\phi}
        }
    \]

    Los números cuánticos cumplen
    \[
        \begin{cases}
            n=1,2,\ldots\\
            l=0,1,\ldots,n-1\\
            m=-l,-l+1,\ldots,l
        \end{cases}
    \]

    Si las funciones angulares están normalizadas:
    \[
        \boxed{
        \int_0^\infty |R_{nl}(r)|^2r^2\,dr=1
        }
    \]

\end{frame}
\begin{frame}{Nivel fundamental ($n=1$)}
 La autoenergía correspondiente es
    \[
        \boxed{E_1=-13.6\ \mathrm{eV}}
    \]

    Para $n=1$, necesariamente $l=0$ y $m=0$.La función radial normalizada es
    \[
        \boxed{
        R_{10}(r)=2a^{-3/2}e^{-r/a}
        }
    \]

    Como $Y_0^0=1/\sqrt{4\pi}$, la función de onda completa resulta
    \[
        \boxed{
        \psi_{100}(\mathbf r)=
        \frac{1}{\sqrt{\pi a^3}}e^{-r/a}
        }
    \]

    \vspace{0.3cm}

   
\end{frame}

\begin{frame}{Segundo nivel de energía ($n=2$)}
    Las soluciones corresponden a la autoenergía:
    \[
        \boxed{E_2=-3.4\ \mathrm{eV}}
    \]
    Para $n=2$, los valores permitidos son $l=0,1$.Las funciones radiales normalizadas son

    \[
    \boxed{
    R_{20}(r)=
    \frac{a^{-3/2}}{\sqrt{2}}
    \left(1-\frac{r}{2a}\right)e^{-r/(2a)}
    }
    \]

    \[
    \boxed{
    R_{21}(r)=
    \frac{a^{-3/2}}{2\sqrt{6}}
    \left(\frac{r}{a}\right)e^{-r/(2a)}
    }
    \]

    \vspace{0.3cm}
\end{frame}
